\SourceLeaf{277}%
substitution of \(90^0-P\), \(90^0-Q\) for \(w,U\), and after applying the transformation given in Article 6 to \(\alpha\cos U\) and \(\alpha\sin U\), while writing \(c,s\) in place of \(\cos P,\sin P\) for short. This gives
\[
\begin{aligned}
p={}&\frac{\cos\theta}{\cos\varphi}\mathord{.}q\\
&-\frac{3ee}{2\cos\varphi^2}\mathord{.}csqq\\
&+\frac{ee}{2\cos\varphi^3\cos\theta}\{-cc+ss+ee(5ccss-s^4)\}q^3\\
&+\frac{ee}{24\cos\varphi^4\cos\theta^2}cs\{16+ee(41cc-77ss)-e^4(101ccss-61s^4)\}q^4
\end{aligned}
\]
\[
\begin{aligned}
&+\frac{ee}{120\cos\varphi^5\cos\theta^3}\{16cc-12ss+ee(41c^4-522ccss+81s^4)\\
&\quad-e^4(538c^4ss-1536ccs^4+126s^6)+e^6(857c^4s^4-1030ccs^6+57s^8)\}q^5\\
&+\text{and so on}
\end{aligned}
\]
\indent The numerical values for our example follow in a form similar to that above, with \(p\) expressed in seconds and \(q\) in degrees.
\[
\begin{aligned}
p={}&360443''852122(\frac{q}{100})\\
&-3052{,}649780(\frac{q}{100})^2\\
&+1002{,}642506(\frac{q}{100})^3\\
&+4119{,}589282(\frac{q}{100})^4\\
&-431{,}181623(\frac{q}{100})^5\quad\text{and so on}
\end{aligned}
\]
\Article{9}
\indent Similarly, the hyperbolic logarithm of \(m\) is expanded as the following series in successive powers of \(q\). The coefficient of \(q^6\) was derived by a different method, with the derivation limited to its principal part.
\[
\begin{aligned}
\log m={}&-\frac{2ee}{3\cos\varphi\cos\theta}\mathord{.}csq^3\\
&-\frac{ee}{6\cos\varphi^2\cos\theta^2}\mathord{.}cc(1-7eess)q^4\\
&+\frac{ee}{30\cos\varphi^3\cos\theta^3}\mathord{.}\frac{s}{c}\{2cc-3ss+ee(20c^4-10ccss+6s^4)\\
&\qquad-e^4(59c^4ss-8ccs^4+3s^6)\}q^5\\
&+\frac{ee}{180\cos\varphi^4\cos\theta^4}\mathord{.}\frac{1}{cc}(2c^4-18ccss-15s^4)q^6
\end{aligned}
\]
